\ArithmeticHead{352}{L}{ANALYSIS.\quad NACHLASS.}
\NumberedItem{9}
\noindent\hspace*{\parindent}Es sei ferner wie oben [art. 7.]
\[
P=F(\alpha,\beta,\gamma,x),\qquad Q=F(\alpha,\beta,\alpha+\beta+1-\gamma,1-x),
\]
\noindent so dass
\[
Q\frac{dP}{dx}-P\frac{dQ}{dx}=\frac{A}{x^\gamma(1-x)^{\alpha+\beta+1-\gamma}},
\]
\noindent wo
\[
A=\frac{\Pi(\gamma-1)\Pi(\alpha+\beta-\gamma)}{\Pi(\alpha-1)\Pi(\beta-1)}\mathpunct{\raisebox{.35ex}{.}}
\]
\noindent Ferner
\[
Q=B{.}\,P-\frac{A}{1-\gamma}{.}\,x^{1-\gamma}F(\alpha+1-\gamma,\beta+1-\gamma,2-\gamma,x),
\]
\noindent wo
\[
B=\frac{\Pi(-\gamma)\Pi(\alpha+\beta-\gamma)}{\Pi(\alpha-\gamma)\Pi(\beta-\gamma)}\mathpunct{\raisebox{.35ex}{.}}
\]
\noindent{\small[S. 68]}\par
\indent Hiedurch kann also die bei $x>0{,}5$ langsamere Convergenz auf eine schnellere gebracht werden. Diess Verfahren ist jedoch nicht anwendbar, wenn $\gamma=1$. Setzt man in diesem Fall
\[
Q=R-AP\log x,
\]
\noindent wo
\[
A=\frac{\Pi(\alpha+\beta-1)}{\Pi(\alpha-1)\Pi(\beta-1)},
\]
\noindent so wird
\[
\begin{aligned}
&R=A\left(-2\Pi'0-\frac{\Pi'(\alpha-1)}{\Pi(\alpha-1)}-\frac{\Pi'(\beta-1)}{\Pi(\beta-1)}\right)+(1)\alpha\beta x\\[.7em]
&\qquad+((1)+(2))\frac{\alpha{.}\,\alpha+1{.}\,\beta{.}\,\beta+1}{1{.}\,2{.}\,1{.}\,2}xx\\[.7em]
&\qquad+((1)+(2)+(3))\frac{\alpha{.}\,\alpha+1{.}\,\alpha+2{.}\,\beta{.}\,\beta+1{.}\,\beta+2}{1{.}\,2{.}\,3{.}\,1{.}\,2{.}\,3}x^3+\mathrm{etc.},
\end{aligned}
\]
\noindent wo
\[
\begin{aligned}
(1)&=\frac{1-(2\alpha-1)(2\beta-1)}{2\alpha\beta}=\frac1\alpha+\frac1\beta-2,\\[.8em]
(2)&=\frac{4-2\alpha{.}\,2\beta}{4{.}\,\alpha+1{.}\,\beta+1}=\frac1{\alpha+1}+\frac1{\beta+1}-\frac22,\\[.8em]
(3)&=\frac{9-(2\alpha+1)(2\beta+1)}{6{.}\,\alpha+2{.}\,\beta+2}=\frac1{\alpha+2}+\frac1{\beta+2}-\frac23
\end{aligned}
\]
\begin{center}etc.\end{center}
\clearpage
